\subsection{Espera e simulação por faixa horária}\label{sec:F-ESPERA}

A estimativa usa as partidas distintas de cada padrão, identificadas pelo
primeiro horário de cada viagem. Aqui, a estimativa por linha é feita por
padrão (sentido e sequência de paradas), sem reunir variantes distintas, dentro de uma faixa $[a,b)$. Após ordenar os
horários $t_1,\ldots,t_n$, calcula-se o intervalo médio
$H=(t_n-t_1)/(n-1)$ e a espera $H/2$. As frequências já foram expandidas pelo
leitor. Padrões com apenas uma partida usam a espera provisória de 300~s,
identificada na interface; padrões sem partidas na faixa não permitem embarque.

A simulação seleciona a faixa pela saída, aplica as esperas às arestas de
embarque de $G_E$ e reutiliza o Dijkstra e a reconstrução da F-ROTA. O resultado
separa caminhada, espera e viagem. A Figura~\ref{fig:atividade-espera}
resume o fluxo. Para $N$ viagens, $P$ padrões e $F$ faixas, a estimativa custa
$O(PF+FN+N\log N)$ de tempo e $O(N+PF)$ de espaço; a consulta acrescenta a
construção de $G_E$ e o Dijkstra, $O((V+E)\log V)$.

\begin{figure}[ht]
  \centering
  \includegraphics[width=0.72\linewidth]{atividade-espera.png}
  \caption{Estimativa de espera e simulação com uma faixa fixa.}
  \label{fig:atividade-espera}
\end{figure}

No instantâneo de 05/11/2026 (SHA-256 iniciado em \texttt{623bdf422fb0}),
serviço de 04/11, o pico da manhã teve 481 padrões com estimativa e 20 com espera
provisória. O grafo teve $V=46.100$ e $E=208.956$. Em macOS arm64 com
Python~3.12.13, cinco repetições de \texttt{medir} produziram média de
5,309~ms para estimar intervalos e 12,441~ms para consultar o grafo pronto.
A função completa de simulação, chamada pela página a cada consulta, levou
248,7~ms em média, incluindo a estimativa, a construção de $G_E$ e a busca.
Essa medição não inclui a renderização da interface nem a leitura do GTFS.
A consulta 14788851 $\rightarrow$ 14792190 às 08h resultou em 40,24~min,
sendo 9,72 de caminhada, 15,34 de espera e 15,18 de viagem. Não é um pior caso.
Comando, desvios e resultados completos constam na feature e no JSON gerado por
\texttt{benchmarks/medir\_espera.py}.

A espera $H/2$ supõe chegadas ao acaso e intervalos regulares. Com intervalos
irregulares, $E[H^2]/(2E[H])$ pode ser maior; essa correção não é aplicada nesta
versão. A faixa da saída permanece fixa durante todo o percurso, mesmo quando
ele atravessa outro período. A Seção~\ref{sec:F-ESPERA-2} apresenta a consulta às partidas programadas em $G_T$.
